\chapter{Background}
\label{cha:background}

\section{Observability and Operational Telemetry}
\label{sec:background:observability}

Observability concerns the extent to which a system's internal behaviour can
be understood through the information it exposes. In distributed software,
this information includes operational telemetry such as logs, metrics and
traces. Monitoring uses such information to track selected aspects of system
behaviour, while observability supports investigation of the conditions
behind that behaviour \parencite{usman2022observability}. These activities
are closely related. For example, an alert about an increased error rate
identifies a condition requiring attention. Explaining the increase requires
evidence about the affected requests and components.

\noindent\textbf{Logs.} A log records events emitted by an application or
another system component. Log entries commonly include timestamps and
messages describing what occurred. A log entry can provide details about a failed operation, but
it is not necessarily associated with a particular user request. Additional
context is needed to connect that event to activity elsewhere in the system
\parencite{opentelemetry2026observability}. This distinction matters when
several requests are processed concurrently: observing an error during an
incident does not, by itself, establish which request it concerns.

\noindent\textbf{Metrics.} Numerical measurements describe system behaviour
over time, such as resource utilisation, request rates and response times.
They support
identification of changes in behaviour and the definition of alert conditions
\parencite{usman2022observability}. Metrics can summarise many observations.
For example, an error rate describes the proportion of unsuccessful requests
over a time interval. That summary can indicate the scale of a problem, while
further information is needed to explain an individual failure.

\noindent\textbf{Distributed traces.} A distributed trace links the recorded
steps involved in handling a request. It is made up of \emph{spans}, each
describing one operation and its timing. For example, a request to a web
application may lead to a database query. That query can be recorded as a
span within the request's trace
\parencite{opentelemetry2026observability}. Linking these spans helps show
which operations belong to the request and where time was spent
\parencite{ashok2024traceweaver}.

To connect operations across services, the services need to pass tracing
information between them. The W3C Trace Context standard defines a common
format for this information, including an identifier for the trace
\parencite{w3c2021tracecontext}.

A trace may still have gaps. If a component is not set up to record tracing
data, its operations may be missing from the trace
\parencite{ashok2024traceweaver}. A missing span therefore does not prove
that the request never reached that component.

The complementary roles of these signals can be illustrated by a hypothetical
application failure. A metric might show an increase in unsuccessful
requests during a particular interval. A trace of one failed request might
identify an unsuccessful database operation, while a log associated with
that operation might record an authentication rejection. The metric
describes a broader change in service behaviour. The linked trace and log
provide evidence about an individual failure. This evidence supports
investigating the authentication rejection, but further information is
needed to explain why it occurred. The rejection alone does not establish
malicious activity or show that other failed requests have the same cause.

Interpreting these observations also requires attention to dependencies.
Research on microservice diagnosis shows how performance problems can
propagate between dependent components, so the location where a symptom is
observed may differ from the origin of the problem
\parencite{gan2021sage}. Correlation helps establish which observations
belong to the same activity, while diagnosis requires assessing what those
observations imply about the failure.

Telemetry can also provide the basis for derived features used in analysis.
In the hypothetical example, a count of recorded authentication rejections
for a service within a defined interval could be calculated from its logs.
This summary describes a pattern across events whose relevance depends on
the problem being investigated. Research on machine learning in security
shows how models can exploit incidental patterns unrelated to their
intended task \parencite{arp2022dosdonts}. This motivates examining the
meaning and diagnostic relevance of telemetry and derived features alongside
their availability and correlation across systems.
